% TOP_PROBLEM: {"problemId": "problem.greek-letter-redshift-conjecture-in-algebraic-k-theory", "canonicalTitle": "Greek-Letter Redshift Conjecture in Algebraic K-Theory", "releaseRank": 453, "primaryDomain": "algebra_representation_category", "primaryDomainLabel": "Algebra, representation & category theory", "displayStatus": "open_with_solved_subcases", "releaseStatus": "open_with_solved_subcases"}
% !TeX program = lualatex
% Definition and sources only; this file is not an LLM solution attempt.
\documentclass[11pt]{article}
\usepackage[margin=25mm]{geometry}
\usepackage{fontspec}
\setmainfont{Latin Modern Roman}
\usepackage{amsmath,amssymb,mathtools}
\usepackage{unicode-math}
\setmathfont{Latin Modern Math}
\usepackage{xurl,hyperref}
\hypersetup{colorlinks=true,urlcolor=blue}
\setcounter{secnumdepth}{0}
\setlength{\parindent}{0pt}
\setlength{\parskip}{5pt}
\setlength{\emergencystretch}{3em}
\begin{document}
\section*{\texorpdfstring{Greek-Letter Redshift Conjecture in Algebraic K-Theory}{Greek-Letter Redshift Conjecture in Algebraic K-Theory}}\label{453}
\subsection{Definitions and mathematical statement}
Fix a prime at which the $n$-th and $(n+1)$-st Greek-letter families in the source exist as nonzero stable homotopy classes. Let $S$ be the sphere spectrum, so $\pi_jS$ is its $j$-th stable homotopy group. Greek-letter elements are the periodic families constructed from powers of chromatic periodicity classes and connecting maps, or equivalently the specified permanent cycles in the chromatic and $BP$-Adams spectral sequences; use Section 1 of Angelini-Knoll for the indexing.

A commutative ring spectrum $R$ detects a family $\{\alpha_k^{(n)}\}$ when every member remains nonzero under the unit:
\[
 \pi_*S\longrightarrow\pi_*R.
\]
Writing $K(R)$ for the algebraic K-theory spectrum of perfect $R$-modules, Conjecture 1.1 asks that detection of the $n$-th family by $R$ imply detection of the $(n+1)$-st by $K(R)$. The source's existence qualification on the families matters, especially at small primes. Detecting a single element or merely proving a height bound is not the entire assertion.
\par\smallskip\textit{Formulation and concept references:} \href{https://arxiv.org/abs/1810.10088}{[S1]}.

\subsection{Short English statement}
Does taking algebraic K-theory raise by one the Greek-letter family detected by a commutative ring spectrum?
\subsection{Sources}
\begin{itemize}
\item[S1]G. Angelini-Knoll, Detecting beta elements in iterated algebraic K-theory, Trans. AMS.
\url{https://arxiv.org/abs/1810.10088}.
\textit{Formulation source. Consulted 24 September 2026: Conjecture 1.1 and the family-detection definition in its full-text HTML version.}
\end{itemize}
\end{document}
